\section{Exception Handling}

\subsection{Section Overview}
Even in the best-written programs, errors can occur at runtime. These errors, or ``exceptions,'' can be caused by invalid user input, hardware failures, network issues, or logic errors. C++ provides a formal mechanism for handling these exceptional circumstances, allowing us to write more robust and fault-tolerant code. In this section, we'll explore the \texttt{try}, \texttt{catch}, and \texttt{throw} keywords.

\subsection{What is Exception Handling?}
Exception handling is a mechanism that allows you to detect and respond to runtime errors in a structured way. Instead of crashing the program, you can ``throw'' an exception when an error occurs and ``catch'' it in a part of the program that is equipped to handle it.

\subsection{Basic Concepts}
\begin{itemize}
    \item \textbf{\texttt{throw}}: The \texttt{throw} keyword is used to signal that an exceptional event has occurred. You can throw values of any type, including integers, strings, or custom objects.
    \item \textbf{\texttt{try}}: The \texttt{try} keyword marks a block of code that may throw an exception.
    \item \textbf{\texttt{catch}}: The \texttt{catch} keyword marks a block of code that handles an exception. The \texttt{catch} block is executed only if an exception of the specified type is thrown from within the corresponding \texttt{try} block.
\end{itemize}

\subsection{Terminating the Program with std::terminate}
If an exception is thrown but not caught anywhere in the program, the program will call \texttt{std::terminate}, which by default aborts the program.

\subsection{Throwing an Exception from a Function}
It is common to throw exceptions from functions when an error condition is met, such as invalid arguments or an inability to perform a task.
\begin{lstlisting}
double divide(int a, int b) {
	if (b == 0) {
		throw std::string("Division by zero error!"); // works, but not ideal
	}
	return static_cast<double>(a) / b;
}
\end{lstlisting}

\begin{remark}
	Throwing a raw \texttt{std::string} (or an \texttt{int}, or any primitive) is legal but non-idiomatic. It means callers can only catch it with \texttt{catch (const std::string \&)} and it cannot be caught by a \texttt{catch (const std::exception \&)} handler. Prefer throwing objects that inherit from \texttt{std::exception}, as shown later in this section.
\end{remark}

\subsection{Handling Exceptions --- Syntax}
You use a \texttt{try...catch} block to handle potential exceptions.
\begin{lstlisting}
int main() {
	try {
		double result = divide(10, 0);
		std::cout << "Result: " << result << std::endl;
	} catch (const std::string &ex) {
		std::cerr << "Error: " << ex << std::endl;
	}
	return 0;
}
\end{lstlisting}

\subsection{Catching Multiple Exceptions}
You can have multiple \texttt{catch} blocks to handle different types of exceptions. The order of the \texttt{catch} blocks matters; they are checked in sequence. You can use \texttt{catch(...)} to catch any type of exception (a ``catch-all'' handler).

\begin{remark}
	Always catch \textbf{more derived} exception types \textit{before} more base types. If you catch \texttt{std::exception} first, it will swallow all derived exceptions and the specialized handlers will never be reached:
	\begin{lstlisting}
		// WRONG order — std::out_of_range is never caught specifically
		} catch (const std::exception &e) { ... }
		} catch (const std::out_of_range &e) { ... } // dead code

		// CORRECT order
		} catch (const std::out_of_range &e) { ... } // caught first
		} catch (const std::exception &e) { ... }    // fallback
	\end{lstlisting}
\end{remark}
\begin{lstlisting}
try {
	// ... code that might throw
} catch (int e) {
	std::cerr << "Caught an integer exception: " << e << std::endl;
} catch (const std::string &e) {
	std::cerr << "Caught a string exception: " << e << std::endl;
} catch (...) {
	std::cerr << "Caught an unknown exception." << std::endl;
}
\end{lstlisting}

\subsection{Stack Unwinding}
When an exception is thrown, the function call stack is ``unwound.'' The program looks for a \texttt{catch} handler in the current function. If one isn't found, the function exits, its local variables are destroyed, and the program looks for a handler in the calling function. This process continues up the stack until a handler is found or the program terminates. This process ensures that destructors of stack-based objects (like smart pointers) are called, which is key to RAII.

\subsection{User-defined Exception Classes}
It is good practice to define your own exception classes, typically inheriting from \texttt{std::exception}. This allows you to convey more detailed information about the error.
\begin{lstlisting}
#include <exception>

class DivideByZeroException : public std::exception {
public:
	// noexcept: this function promises never to throw an exception itself
	const char* what() const noexcept override {
		return "Attempted to divide by zero";
	}
};

double divide(int a, int b) {
	if (b == 0) {
		throw DivideByZeroException();
	}
	return static_cast<double>(a) / b;
}
\end{lstlisting}

\subsection{The Standard Library Exception Hierarchy}
The C++ standard library provides a hierarchy of exception classes in \texttt{<stdexcept>}, all derived from \texttt{std::exception}. Using these makes your code understandable to anyone familiar with C++.

\begin{lstlisting}
#include <stdexcept>

// std::invalid_argument — bad parameter value
double square_root(double x) {
	if (x < 0)
		throw std::invalid_argument("square_root requires non-negative input");
	return std::sqrt(x);
}

// std::out_of_range — index or value outside valid range
int get_element(const std::vector<int> &v, int i) {
	if (i < 0 || i >= static_cast<int>(v.size()))
		throw std::out_of_range("index out of bounds");
	return v[i];
}

// std::runtime_error — errors only detectable at runtime
void load_file(const std::string &path) {
	// if file not found:
	throw std::runtime_error("Cannot open file: " + path);
}
\end{lstlisting}

\begin{remark}
	C++ has no \texttt{finally} block (unlike Java or Python). Instead, C++ uses \textbf{RAII}: resources are tied to object lifetimes, so their destructors run automatically during stack unwinding even when an exception is thrown. Smart pointers, file stream objects, and mutex locks all use this pattern — no \texttt{finally} needed.
\end{remark}

\subsection{Section Challenge}
Create a simple \texttt{Account} class with a balance. Implement a \texttt{withdraw} method that throws an \texttt{InsufficientFundsException} (a custom exception class you create) if the withdrawal amount is greater than the balance. In \texttt{main}, create an account and wrap a withdrawal attempt in a \texttt{try...catch} block to handle the exception gracefully.

\subsection{Section Challenge --- Solution}
\begin{lstlisting}
#include <iostream>
#include <exception>

class InsufficientFundsException : public std::exception {
public:
	const char* what() const noexcept override {
		return "Insufficient funds for withdrawal.";
	}
};

class Account {
private:
	double balance;
public:
	Account(double bal) : balance(bal) {}
	void withdraw(double amount) {
		if (amount > balance) {
			throw InsufficientFundsException();
		}
		balance -= amount;
		std::cout << "Withdrawal successful. Remaining balance: " << balance << std::endl;
	}
};

int main() {
	Account my_account(100.0);
	try {
		my_account.withdraw(50.0);  // This should succeed
		my_account.withdraw(70.0);  // This should fail and throw
	} catch (const InsufficientFundsException &ex) {
		std::cerr << "Error: " << ex.what() << std::endl;
	}
	return 0;
}
\end{lstlisting}

\subsection{Quiz 15: Section 18 Quiz}
\begin{enumerate}
    \item What happens if an exception is thrown but never caught?
    \begin{enumerate}
        \item[a)] The program continues as if nothing happened.
        \item[b)] The exception is ignored.
        \item[c)] The program calls \texttt{std::terminate}, which usually aborts the program.
        \item[d)] The compiler issues a warning.
    \end{enumerate}

    \item What is the purpose of a \texttt{try} block?
    \begin{enumerate}
        \item[a)] To declare a new type of exception.
        \item[b)] To handle an exception.
        \item[c)] To enclose code that might throw an exception.
        \item[d)] To throw an exception.
    \end{enumerate}

    \item Why is it beneficial to create custom exception classes that inherit from \texttt{std::exception}?
    \begin{enumerate}
        \item[a)] It's the only way to throw an exception.
        \item[b)] It allows your exceptions to be caught by handlers for \texttt{std::exception}, promoting a standard way of handling errors.
        \item[c)] It makes the program run faster.
        \item[d)] It's required for \texttt{try...catch} to work.
    \end{enumerate}
\end{enumerate}

\textbf{Answers:} 1. (c), 2. (c), 3. (b)
